\section{Conclusion}\label{sec:conclusion}

We asked what a large cold multiple-choice lift from contrastive-correctness PEFT
is, on a frontier 31B instruction model, and the answer is that it is not a
capability gain. A capacity- and budget-matched non-contrastive control attributes
most of the lift --- ${\sim}62\%$ of the ARC gain --- to ordinary fine-tuning
rather than to the contrastive objective. The recovery lives in a learned direction
that is off-axis of degree to the model's own correctness geometry, and that
geometry is causal and identifiable on both the harness and the write sides. And
the lift does not become generation: the items it fixes already sit at the base's
chain-of-thought ceiling, a read-only selector on the same activations recovers it
without the deployment costs, and --- most plainly --- the base asked to emit the
answer letter in a single forward recovers more of the base-wrong items ($0.970$)
than the adapter ($0.763$) at the same cost, tying a reasoning baseline that costs
a thousand times more. The effect and its off-axis geometry replicate on a second
model family --- though only partly: the cross-family lift is same-sign on all four
tasks but marginal on Winogrande at $n{=}400$, and the causal spine that carries
the attribution (the arbiter, the inducibility control, the recovery spectrum)
remains single-model.

The practical consequence is a caution for the contrastive-PEFT-for-multiple-choice
program. A cold multiple-choice gain of this size can be an off-axis,
non-transferring edit to a lossy readout --- cheaper to obtain with a read-only
selector, and beaten by a reasoning baseline --- so a benchmark improvement of this
kind is not, on its own, evidence of a better model.
